%======================================================================
\chapter{Introduction}
%======================================================================

The transition to \glsxtrfull{ev} represents a fundamental shift in transportation infrastructure requirements across Canada. Recent data from the Canada Energy Regulator (2024) indicates unprecedented growth in \glsxtrshort{ev} adoption, with sales reaching 12.8\% of new vehicle registrations in 2023, marking a significant milestone in sustainable transportation \cite{cer_market_2024}. This rapid adoption aligns with federal initiatives supporting the Paris Agreement and stems from policy frameworks such as the Canadian Net Zero Emissions Accountability Act, which mandates complete transition to zero-emission vehicle sales by 2035 \cite{ecc_what_2024}.

The \glsxtrfull{kwc-cma} exemplifies both the opportunities and challenges of this transition. Despite strong \glsxtrshort{ev} adoption trends documented in \gls{odc} \cite{ontario_data_2024}, our analysis reveals that current charging infrastructure reaches only 14.57\% of the population for \gls{l2} charging and a mere 1.78\% for \gls{l3} charging, creating a critical disparity between \glsxtrshort{ev} adoption and infrastructure availability. The Hydro One Kitchener-Waterloo-Cambridge-Guelph Regional Infrastructure Plan (2021) highlights the importance of coordinated infrastructure development to support increasing electrical demand from \glsxtrshort{ev} adoption \cite{hydro_kwcg_2021} \cite{hydro_kwcg_2024}.

This study addresses these challenges through a comprehensive optimization framework to improve the \glsxtrshort{kwc-cma} \glsxtrshort{ev} charging infrastructure. To contextualize this framework, it is essential to understand the technological landscape of \glsxtrshort{ev} charging equipment, which is typically classified into three distinct tiers. \gls{l1} chargers utilize standard 120V \glsxtrfull{ac} outlets, delivering approximately 1.4 to 1.9 kW of power. While highly accessible, L1 charging is characterized by prolonged charging durations, often requiring overnight sessions to meaningfully replenish a vehicle's range, making it primarily suited for residential applications. \gls{l2} charging represents a significant step up, operating on 208V to 240V \gls{ac} supplies and providing between 3.2 kW and 19.2 kW of power. Literature indicates that L2 infrastructure serves as the backbone of public and workplace charging networks, offering a pragmatic balance between installation costs and moderate charging speeds that can service vehicles over several hours. In contrast, \gls{l3}, or \glsxtrfull{dcfc}, supplies high-voltage direct current directly to the battery, bypassing the onboard converter. \gls{l3} stations offer vastly superior power supply capacities, typically ranging from 50 kW to over 350 kW, enabling them to recharge an \glsxtrshort{ev} battery to 80\% capacity in under 30 minutes. This rapid charging speed is critical for facilitating long-distance travel and alleviating range anxiety among drivers. We employ a \glsxtrlong{mip} to optimize the placement of new charging stations while considering strategic upgrades of existing \gls{l2} stations to \gls{l3} capabilities, directly addressing the varying power demands and duration constraints identified in recent literature. Our approach integrates multiple data streams to develop cost-effective solutions that maximize charging coverage. The insights gained from this study can inform policy development, guide infrastructure investments, and support the broader transition to sustainable urban mobility.
